% !TeX program = xelatex
\documentclass[aspectratio=169,11pt]{beamer}

% Put beamerthemeMyPre.sty in the same directory as this file.
\usetheme{MyPre}

% -----------------------------------------------------------------------------
% Basic information
% -----------------------------------------------------------------------------
\title{Function Secret Sharing:\\[-1mm]Distributed Comparison Function}
\author{Artemis}
\date{2026.7.15}

\begin{document}

% Cover: deliberately unnumbered.
\begin{frame}[plain,noframenumbering]
  \titlepage
\end{frame}

% -----------------------------------------------------------------------------
% Contents page
% It appears before the first numbered section, so the theme renders "Contents"
% as a single first-level title and keeps the following section numbering at I.
% -----------------------------------------------------------------------------
\begin{frame}[t]{Contents}
  % Demo: vertical top alignment + horizontal left alignment.
  \begin{MyPreContent}[top][left]
    \vspace{-1.5mm}
    \tableofcontents
  \end{MyPreContent}
\end{frame}

% -----------------------------------------------------------------------------
% Manual alignment for each content group
% Keep the frame itself as [t], then wrap a group with MyPreContent:
%   \begin{MyPreContent}[top][left]       ... \end{MyPreContent}
%   \begin{MyPreContent}[top][center]     ... \end{MyPreContent}
%   \begin{MyPreContent}[center][left]    ... \end{MyPreContent}
%   \begin{MyPreContent}[center][center]  ... \end{MyPreContent}
% -----------------------------------------------------------------------------

% -----------------------------------------------------------------------------
% Example section + frames
% Level 1 title: I. Review
% Level 2 title: only the frame topic, e.g. DPF Invariants
% -----------------------------------------------------------------------------
\section{Review}

% Each second-level heading is declared as a subsection so it also appears in
% the contents page.  The frame title uses the same text, which keeps the
% upper-left two-level heading visually consistent.
\subsection{Distributed Point Function}
\begin{frame}[t]{Distributed Point Function}
  \begin{MyPreContent}[top][left]
  \begin{block}{Definition}
    A point function can be written as
    \[
      f_{\alpha,\beta}(x)=
      \begin{cases}
        \beta, & x=\alpha,\\
        0,     & x\neq\alpha.
      \end{cases}
    \]
  \end{block}

  \vspace{2mm}
  \begin{itemize}
    \item This area intentionally uses a conventional mathematical Beamer style.
    \item English text is Times New Roman; Chinese text is 楷体。
    \item The page number is fixed at the lower-right corner.
  \end{itemize}
  \end{MyPreContent}
\end{frame}

\subsection{DPF Invariants}
\begin{frame}[t]{DPF Invariants}
  \begin{MyPreContent}[top][left]
  Let $k_b$ be the key held by party $P_b$. A typical invariant is
  \[
    t_0 \oplus t_1 =
    \begin{cases}
      1, & \text{on the special path},\\
      0, & \text{otherwise}.
    \end{cases}
  \]

  \begin{theorem}[Example]
    If the correction words preserve the invariant at every level, the two
    evaluations reconstruct the desired point-function output.
  \end{theorem}

  \begin{proof}
    Apply the invariant recursively from the root to the leaf, then combine the
    final correction word with the leaf payload.
  \end{proof}
  \end{MyPreContent}
\end{frame}

\section{Construction}

\subsection{A Classical Mathematical Slide}
\begin{frame}[t]{A Classical Mathematical Slide}
  \begin{MyPreContent}[top][left]
  \begin{columns}[T,onlytextwidth]
    \column{0.49\textwidth}
      \begin{block}{Protocol sketch}
        \begin{enumerate}
          \item Sample seeds.
          \item Expand with a PRG.
          \item Apply correction words.
          \item Reconstruct the output.
        \end{enumerate}
      \end{block}

    \column{0.47\textwidth}
      \[
        \Pr[\mathsf{View}_0 \approx_c \mathsf{View}_1]
        \ge 1-\operatorname{negl}(\lambda).
      \]

      中文示例：这一段会自动使用楷体；英文仍使用 Times New Roman。
  \end{columns}
  \end{MyPreContent}
\end{frame}

\end{document}
